\documentclass[10pt,a4paper]{amsart}
\usepackage[margin=19mm]{geometry}
\usepackage{amsmath,amssymb,mathtools}
\usepackage[scr=rsfs,cal=euler]{mathalfa}
\usepackage{xfrac}
\usepackage[unicode,hidelinks,bookmarks=false]{hyperref}
\newcommand{\ie}{i.e.\ }
\newcommand{\cf}{cf.\ }
\newcommand{\resp}{respectively\ }
\newcommand{\Subsection}{\subsection}
\providecommand{\nobreakdash}{\nobreak\hbox{-}}
\setlength{\emergencystretch}{2em}
\multlinegap0pt
\usepackage{bm}
\usepackage{bbm}
\usepackage{dsfont}
\usepackage{xspace}
\usepackage{cancel}
\usepackage{mathtools}
\usepackage[shortlabels]{enumitem}
\usepackage{amsthm}
\makeatletter
\@ifundefined{theorem}{\newtheorem{theorem}{Theorem}}{}
\@ifundefined{lemma}{\newtheorem{lemma}[theorem]{Lemma}}{}
\makeatother
\newcommand{\AKone}{B}
\newcommand{\cK}{{\mathcal K}}
\newcommand{\kip}{[\,\cdot\, , \cdot\,]}
\newcommand{\mD}{\mathcal{D}}
\newcommand{\scalar}[2]{\langle#1,\, #2\rangle} 
\title[1-D Schr\"odinger operator on a star graph with nondefinite ]{1-D Schr\"odinger operator on a star graph with nondefinite weight function}
\author{Edison Leguizam\'on, Carsten Trunk, Mitsuru Wilson, Monika Winklmeier}
\date{Source: arXiv:2505.22901v2 (CC BY 4.0)}
\begin{document}
\maketitle

\noindent\emph{Source and licence.} 1-D Schr\"odinger operator on a star graph with nondefinite weight function. Version arXiv:2505.22901v2 (CC BY 4.0), distributed under the Creative Commons Attribution 4.0 International licence, as recorded in the arXiv licence metadata for this exact version and shown on the abstract page https://arxiv.org/abs/2505.22901v2. Full rights notice: licenses/arxiv-2505.22901-CC-BY-4.0.txt.\par\smallskip

\noindent\textit{Editorial scope.} Counted unit: the Theorem whose source label is \texttt{None} in the article (source0.tex lines 1064--1066, source label \texttt{None}), delivered together with the article's own complete original proof of it (source0.tex lines 1067--1112, one \texttt{proof} environment of 46 lines). No mathematics is written, repaired or re-derived: statement and proof are the article's own text line for line. Required context retained uncounted: Halle (lines 269--272); Goettingen (lines 298--318); Goettingen:i (lines 302--311); Goettingen:ii (lines 305--307); lemma:invertible\_operator (lines 966--977); item:Xi (lines 970--976); item:Xii (lines 970--976); lemma:form\_domain (lines 1018--1025). Every definition, display and label of those lines is retained unchanged. Omitted from this package and not redistributed: 1--268, 273--297, 308--965, 977--1017, 1026--1063, 1113--1134. Nothing inside a retained excerpt is omitted or altered: every retained body is delivered exactly as the article's lines are written, so no omission receipt of any kind accompanies this package. Citations: the retained excerpts carry no citation command at all, so no bibliography entry of the article is transcribed and no reference is redistributed. Reference adaptation: none was needed and none was made; every reference inside the delivered unit resolves inside it, so no unresolved reference or citation is produced. Printed slips kept literal: no printed word, symbol, hypothesis or number of the counted statement or of its proof was altered. Layout and class: the article's own class is \texttt{article} with packages including amsmath, amsthm, amssymb, enumitem, xcolor, geometry, mathrsfs, titling, todonotes, tcolorbox, tikz, tikzrput, pgfplots, fancyhdr; the wrapper is the lane's pinned \texttt{amsart} editorial wrapper at 10pt on A4, which restates the theorem-like environments the retained text needs and adds the article's own macro definitions that the retained text needs (\texttt{\textbackslash AKone}, \texttt{\textbackslash cK}, \texttt{\textbackslash kip}, \texttt{\textbackslash mD}, \texttt{\textbackslash scalar}), with extra wrapper packages (bm, bbm, dsfont, xspace, cancel, mathtools, enumitem). The delivered numbering is the unit's own. Assets: no retained span contains a figure environment, a graphics-inclusion command, a PDF-page inclusion, a table or a TikZ picture; no image file is copied into this package, the package declares no asset dependency and no image file is redistributed with it. Licence: version arXiv:2505.22901v2 of this article is distributed under the Creative Commons Attribution 4.0 International licence, as recorded in the arXiv licence metadata for this exact version and shown on the abstract page https://arxiv.org/abs/2505.22901v2. A full rights notice is delivered at licenses/arxiv-2505.22901-CC-BY-4.0.txt. Characters: every retained excerpt is pure ASCII, so the delivered engine, XeLaTeX with its default Latin Modern text font, reports no missing character.
\begin{equation}
\label{Halle}
	\scalar{x}{y} = [Jx,y], \quad  x,y \in \cK,
\end{equation}

\begin{theorem}
\label{Goettingen}
Let $A$ be a non-negative operator in a Krein space $\bigl(\cK,\kip\bigr)$ with $0\in\rho(A)$. Then, the following two statements are equivalent.

\begin{enumerate}[label={\upshape(\roman*)}]
    \item\label{Goettingen:i}
    There exists a bounded operator $W$, positive with respect to the Krein space inner product $\kip$ with a bounded everywhere defined inverse such that
    \begin{equation*}
	W\mD\left((JA)^{\frac{1}{2}}\right)\subset \mD\left((JA)^{\frac{1}{2}}\right). 
    \end{equation*}
		
    \item\label{Goettingen:ii}
    The operator $A$ is similar to a selfadjoint operator in the Hilbert space $\bigl(\cK,\scalar{\cdot}{\cdot}\bigr)$, where $\scalar{\cdot}{\cdot}$ is given by \eqref{Halle}.
\end{enumerate}
If the spectrum of $A$ is discrete, the above statements \ref{Goettingen:i} and \ref{Goettingen:ii} are equivalent to the following statement.
\begin{enumerate}[label={\upshape(\roman*)}]
\setcounter{enumi}{2}
    \item\label{Goettingen:iii}
    There exists an Riesz basis in $\cK$ consisting of eigenfunctions of $A$.
\end{enumerate}
\end{theorem}

\begin{lemma}
\label{lemma:invertible_operator}
   For every $\delta\in (0, 1/2)$ there exists a positive bounded and boundedly invertible operator $X: L^2(G)\to L^2(G)$ in the Hilbert space $(L^2(G),\scalar{\cdot}{\cdot})$ such that
   for every $f = (f_j)_{j=1}^n\in L^2(G)$ the following is true.
   \begin{enumerate}[label={(\roman*)}]
	\item\label{item:Xi} 
    $(f-Xf)_j(x) = 0$ for $x \in(1-\delta, 1)\subseteq I_j$;
    
	\item\label{item:Xii} 
        if $f_j$ has a continuous extension to $0$, then so does $(Xf)_j$ and $(Xf)_j(0) = 0$.
   \end{enumerate}
\end{lemma}

\begin{lemma}
   \label{lemma:form_domain}
   The form domain $\mD[J\AKone]$ of $J\AKone$ is equal to
   \begin{align}
      \label{eq:formdomainAone}
      \mD[J\AKone] = \left\{ f\in \widehat H^1_D(G)  : f_1(0) = \dots = f_n(0) \right\}.
   \end{align}
\end{lemma}

\begin{theorem}
   The operator $\AKone$ is similar to a selfadjoint operator in the Hilbert space $L^2(G)$. In particular, its eigenfunctions form a Riesz basis.
\end{theorem}

\begin{proof}
   Since $\AKone$ is nonnegative with $0\in \rho(\AKone)$ it remains to show that there exists a positive, bounded and boundedly invertible operator 

   $$
   W:L^2(G) \to L^2(G) \quad \mbox{such that} \quad 
   W(\mD(( J\AKone )^{\frac{1}{2}}))) \subseteq 
   \mD(( J\AKone )^{\frac{1}{2}}),
   $$
   cf.\ Theorem~\ref{Goettingen}.
   We claim that the operator $W = JX$ with $X$ from Lemma~\ref{lemma:invertible_operator} is such an operator.
   Clearly, $X$ is bounded and boundedly invertible in $L^2(G)$, hence $W$ is bounded and boundedly invertible.
   Moreover, the positivity of $X$ in the Hilbert space $(L^2(G),\scalar{\cdot}{\cdot})$ gives the positivity of $W$ in the Krein space $(L^2(G),\kip)$.
   By construction, 
   $$
   Xf\in \bigoplus_{j=1}^n H^1(I_j) \quad 
   \mbox{if} \quad f\in \bigoplus_{j=1}^n H^1(I_j),
   $$
   hence by Lemma \ref{lemma:form_domain}
   $$
   X(\mD[J\AKone])\subseteq X\Big(\bigoplus_{j=1}^n H^1(I_j)\Big)\subseteq \bigoplus_{j=1}^n H^1(I_j).
   $$
   If $f\in \mD( (J\AKone)^{\frac{1}{2}}) )$, then $f_j(1) = 0$ for all $j=1,\dots, n$, 
   hence \ref{item:Xi} from Lemma~\ref{lemma:invertible_operator} shows that 
   $$
   f_1(1) = \dots = f_n(1) = 0
   $$
   and \ref{item:Xii} from the same lemma gives 
   $$
   f_1(0) = \dots = f_n(0) = 0.
   $$
   In summary, 
   $$
   Xf\in \bigoplus_{j=1}^n H_0^1(I_j).
   $$
   Clearly 
   $$
   J \Big( \bigoplus_{j=1}^n H_0^1(I_j) \Big) = \bigoplus_{j=1}^n H_0^1(I_j) \subseteq \mD[J\AKone].
   $$
   Therefore
   \begin{equation*}
      W(\mD[J\AKone]) = JX(\mD[J\AKone])
      \subseteq J \Big( \bigoplus_{j=1}^n H_0^1(I_j) \Big)
      \subseteq \mD[J\AKone].
      \qedhere
   \end{equation*}
\end{proof}

\end{document}
